\scanpage{68}
\begin{center}6 - 33 -\hfill n\textsuperscript{o} 392\end{center}
\begin{center}§ 4.\end{center}

1) Soient $A$ un anneau noethérien intégralement clos, $V$ un espace vectoriel de dimension finie sur le corps des fractions de $A$. Montrer que si $(M_\lambda)$ est une famille quelconque de réseaux réflexifs de $V$ contenant toutes un même réseau $N$, le réseau $M=\bigcap_\lambda M_\lambda$ est réflexif, (considérer les réseaux associés $M_\lambda^*$).

*2) Soient $A$ un anneau noethérien intégralement clos, $E$, $F$ deux $A$-modules de type fini. On suppose que $E$ est sans torsion, que $\operatorname{Hom}_A(E,F)$ est réflexif et que $\operatorname{Ext}_A^1(E,F)=0$. Montrer alors que $F$ est réflexif. (Prouver d'abord que $F$ est sans torsion ; si $T=F^{**}/c_F(F)$, calculer de deux manières $\operatorname{Ass}(\operatorname{Hom}_A(E,T))$ en utilisant la suite exacte :
\diagram{\begin{tikzcd}
0 \arrow[r] & \operatorname{Hom}(E,F) \arrow[r] & \operatorname{Hom}(E,F^{**}) \arrow[r] & \operatorname{Hom}(E,T) \arrow[r] & 0
\end{tikzcd}}
et le chap. IV, § 1, n\textsuperscript{o}\quad, prop.\quad). *

3) Soient $A$ un anneau noethérien intégralement clos, $K$ son corps des fractions, $M$ un réseau réflexif d'un espace vectoriel $V$ de dimension finie sur $K$, $L$ un réseau \uline{libre} de $V$ contenant $M$. Montrer qu'il existe un réseau libre $L_1$ de $V$ tel que $M=L\cap L_1$. (Considérer l'ensemble fini $I$ des idéaux premiers $\underline p$ de hauteur 1 de $A$ tels que $L_{\underline p}\ne M_{\underline p}$, et l'anneau principal $S^{-1}A$, où $S=\bigcap_{\underline p\in I}(A-\underline p)$ ; montrer qu'il existe un réseau libre $L_0$ de $V$ tel que $M_{\underline p}=(L_0)_{\underline p}$ pour tout $\underline p\in I$, et un $s\in S$ tel que $M\subset s^{-1}L_0=L_1$).

4) Soient $k$ un corps, $A$ l'anneau de polynômes $k[X,Y]$.

a) Soit $(e_1,e_2)$ la base canonique du $A$-module $A^2$, et soit $E$ le sous-$A$-module de $A^2$ engendré par $(X-Y)e_1$, $e_1+Xe_2$ et $e_1+Ye_2$ ; soit $F$ le sous-module monogène de $E$ engendré par $(X-Y)^2e_1$. Montrer que le $A$-module $M=E/F$ n'est pas somme directe de son sous-module de torsion et d'un module sans torsion.

b) Montrer que le $A$-module de torsion $A/AXY$ n'est pas somme directe de sous-modules monogènes de la forme $A/\underline P_i^{n_i}$, où les $\underline P_i$ sont des idéaux premiers de hauteur 1.

5) Soient $A$ un anneau noethérien intégralement clos, $M$ un $A$-module de type fini. Montrer que si $\underline a$, $\underline b$ sont deux idéaux de $A$, le $A$-module $((\underline aM)\cap(\underline bM))/(\underline a\cap\underline b)M$ est pseudo-nul ; donner un exemple où il n'est pas nul.

\margin{67}
